\section{Related Works}

This paper is closely related to auto-regressive methods for text conditioned image and video generation. DALL-E~\cite{DALLE} \textit{translates} text tokens to discrete image embeddings learnt using a VQVAE~\cite{vqvae}. Parti~\cite{PARTI} has a similar architecture but can generate higher quality images by predicting tokens from a ViT-VQGAN~\cite{vit_vqgan} using a 21B parameters transformer. Similar architectures have been used for generating videos as well. GODIVA~\cite{godiva} uses a transformer to map text tokens to video tokens from a image based VQVAE. Given the large number of tokens from multiple frames, GODIVA relied on a local-attention mechanism. Similarly, NUWA~\cite{nuwa} and NUWA-Infinity~\cite{wu2022nuwa} both employ auto-regressive architectures to generate videos and images from text. NUWA generates fixed size outputs, while NUWA-Infinity introduces a second layer of auto-regressive computation to support variable size videos. Likewise, CogVideo~\cite{cogvideo} argues the main reason behind low quality video generation is the scarcity of good text-video data and tried to leverage pre-trained text to images models to generate high quality video. 

While \model sticks to the same architecture principles, it has major differences with previous work. Most notably, NUWA, NUWA-Infinity and CogVideo treat videos as a sequence of independent images. This can lead to poor modeling of dynamics and generate motion artifacts. To combat this, NUWA-infinity used the previous frame during decoding to combat this. In \model, we go further and treat videos as a temporal sequence of images which substantially decreases the number of video tokens given the redundancy in video generation, and results in a much lower training cost. The auto-regressive nature of the \phenaki also allows us to effectively condition on previous frames and generates longer videos as detailed in Section~\ref{sec:method}.

%\mbz{I'm not sure how much detail we wanna put here. commenting out for now.}
% Second, \model does not rely on an auto-regressive transformer. 
% \begin{itemize}
%     \item It is similar to maskgit
%     \item It is much faster in inference
%     \item Inference cost can actually be tuned, we can trade off quality for generation speed
%     \item The efficient maskgit inference opens the door to story mode. Which would be harder with regular transformers. 
% \end{itemize}

% Finally, 
% \begin{itemize}
%     \item CogVideo does pre-training
%     \item Nuwa does train on text to image and text2video jointly + a host of other tasks. In our model, text2video and text2image are actual naturally co-existing leading to a less complicated setup.
%     \item similar to imagen and video diffusion we use a pre-trained t5x text encoder. This difffers
% \end{itemize}

%Second, \model is the first model that utilizes joint training on images and videos at the same time to rectify issues with text-video datasets. Finally, \model is only about 1B parameters in size, which is much smaller than CogVideo at 9B and NUWA which only has a 800M parameters decoder. This smaller size, clearly highlights the importance of main contributions of our paper.
% \pikinder{Nuwa Looks to be 870M in total }

Diffusion models are another class of models which recently have been used for conditional and unconditional video generation, which we call VDM~\cite{ho2022video}. In VDM, authors proposed replacing the conventional U-Net architectures for 2D image modeling with a 3D space-time model to run the diffusion process directly on pixels. % VDM then uses two models to generate videos: one to generate every 4 frames and one to fills in between. 
While this approach provides an effective formulation for modeling videos, it is limited to fixed size videos. To address this issue, VDM provides an auto-regressive extension, which allows the model to generate longer videos but it is typically impractical due to high sampling time of diffusion models.

Text conditional video generation is a relatively new field of research, nonetheless, image conditional video generation, commonly known as video prediction, and unconditional video generation have been studied more comprehensively. These papers include deterministic methods using a combination of recurrent and convolutional networks~\cite{ranzato2014video,srivastava2015unsupervised,finn2016unsupervised, wang2017predrnn}, variational based stochastic methods~\cite{babaeizadeh2017stochastic,svg,villegas2019high,fitvid} and more recently by learning a discrete representation~\cite{videovqvae,rakhimov2020latent,transframer}, auto-regressive models~\cite{videotranf,videogpt,ccvs,harp}, diffusion models~\cite{mcvd,harvey2022flexible,yang2022diffusion,diff4pred} flow based models~\cite{kumar2019videoflow}, and finally adversarial based methods~\cite{vondrick2016generating,saito2017temporal,tulyakov2018mocogan,clark2019adversarial,saito2020train,trIVD}. These works mostly consider limited domain (e.g. robotic videos) prediction/generation, or short fixed size clips. Section~\ref{sec:exp} provides comparison with some of these models.
